\documentclass[11pt]{article}
\input{style-sujet}
\usepackage{array}

\begin{document}

\entetesujet{Sujet bac 2018 -- Série C}

% ====================== EXERCICE 1 ======================
\exo{1}{4 points}

On considère le système d'équations $(S)$ défini par :
\[ (S) : \begin{cases} x \equiv 2\ [36] \\ x \equiv 3\ [25] \end{cases} \]

\begin{enumerate}
  \item Montrer que le système $(S)$ est équivalent à l'équation $(E) : 36a - 25b = 1$ où $a$ et $b$ désignent des entiers relatifs.
  \item Vérifier que le couple $(-9, -13)$ est une solution de $(E)$.
  \item Montrer que l'équation $(E)$ est équivalente à l'équation $(E') : 36a \equiv 1\ [25]$.
  \item Donner l'inverse modulo 25 de 36.
  \item En déduire les solutions de $(E')$.
  \item \begin{enumerate}[label=\alph*.]
    \item Déterminer les solutions de l'équation $(E)$.
    \item En déduire les solutions du système $(S)$ telles que $0 < x < 50$.
  \end{enumerate}
\end{enumerate}

% ====================== EXERCICE 2 ======================
\exo{2}{8 points}

Dans le plan orienté $(P)$, on considère un carré $ABCD$ de centre $O$ tel que $(\vv{AB},\vv{AD}) = \dfrac{\pi}{2}\ [2\pi]$.

$E$, $F$, $G$ et $H$ désignent les milieux respectifs des segments $[AB]$, $[BC]$, $[CD]$, $[DA]$. $(C_1)$ est le cercle circonscrit au triangle $ABD$. $I$ est le symétrique de $O$ par rapport à $G$.

\begin{enumerate}
  \item Faire une figure que l'on complétera. On prendra $AB = 4$ cm et on placera $(AB)$ horizontalement.
  \item Soit $(\Gamma)$ l'hyperbole de rectangle fondamental le carré $ABCD$ et d'axe non focal la droite $(EG)$.
  \begin{enumerate}[label=\alph*.]
    \item Préciser le foyer $J$ de $(\Gamma)$ situé sur la demi-droite $[OF)$.
    \item Préciser la directrice $(D)$ de $(\Gamma)$ associée au foyer $J$.
    \item Construire le point $K$ de $(\Gamma)$ situé sur le segment $[JB]$.
    \item Déterminer la demi-droite asymptote de $(\Gamma)$ située dans la portion du plan délimitée par les demi-droites $[OF)$ et $[OG)$.
    \item Construire la branche $(\Gamma_0)$ de $(\Gamma)$ située dans la portion du plan délimitée par les demi-droites $[OF)$ et $[OG)$.
  \end{enumerate}
  \item Soit $S$ la similitude plane indirecte d'axe la droite $(AC)$ de rapport 2, qui transforme le point $F$ en le point $I$. Démontrer que son centre est $O$.
  \item Soit $(\Gamma')$ l'image de $(\Gamma)$ par $S$.
  \begin{enumerate}[label=\alph*.]
    \item Montrer que $(\Gamma')$ est une hyperbole équilatère.
    \item Trouver l'excentricité de $(\Gamma')$.
    \item Construire le cercle principal $(C_2)$ de $(\Gamma')$.
    \item Démontrer que $(\Gamma)$ et $(\Gamma')$ ont les mêmes asymptotes.
    \item Déterminer l'axe focal de $(\Gamma')$.
    \item Construire l'image $J'$ du foyer $J$ de $(\Gamma)$.
    \item Construire l'image $(C_1')$ du cercle $(C_1)$ par $S$.
    \item Construire l'image $(\Gamma_0')$ de $(\Gamma_0)$ par $S$.
  \end{enumerate}
\end{enumerate}

% ====================== EXERCICE 3 ======================
\exo{3}{5 points}

\partie{A}
Soit $g$ la fonction numérique définie sur $]0\,;+\infty[$ par : $g(x) = xe^x - 1$.

\begin{enumerate}
  \item Calculer la dérivée $g'$ de $g$.
  \item Dresser le tableau de variation de $g$. On admet que $\displaystyle\lim_{x\to+\infty} g(x) = +\infty$.
  \item \begin{enumerate}[label=\alph*.]
    \item Montrer que l'équation $g(x) = 0$ admet une solution unique $\alpha$ appartenant à l'intervalle $\left]\dfrac{1}{2}\,;1\right[$.
    \item En déduire le signe de $g(x)$ sur $]0\,;+\infty[$.
  \end{enumerate}
\end{enumerate}

\partie{B}
Soit $f$ la fonction définie sur $]0\,;+\infty[$ par $f(x) = e^x - \ln x$. On désigne par $(C)$ sa courbe représentative dans le plan muni d'un repère orthonormé $(O,\vi,\vj)$.

\begin{enumerate}[start=4]
  \item Calculer les limites de $f$ en $0^+$ et $+\infty$.
  \item Montrer que la dérivée $f'$ de $f$ est $f'(x) = \dfrac{g(x)}{x}$.
  \item \begin{enumerate}[label=\alph*.]
    \item Dresser le tableau de variation de $f$. On admet que la courbe $(C)$ admet une branche parabolique de direction $(Oy)$.
    \item Montrer que $f$ admet un minimum $f(\alpha) = \alpha + \dfrac{1}{\alpha}$.
  \end{enumerate}
  \item Tracer la courbe $(C)$. On prendra $\alpha = 0{,}6$ et $f(\alpha) = 2{,}3$.
\end{enumerate}

% ====================== EXERCICE 4 ======================
\exo{4}{3 points}

Soit la série statistique à deux caractères $(X, Y)$ donnée par le tableau à double entrée ci-dessous.

\begin{center}
\renewcommand{\arraystretch}{1.3}
\begin{tabular}{|c|c|c|c|}
\hline
$X\,\backslash\,Y$ & -1 & 2 & 3 \\
\hline
2 & 2 & 0 & 3 \\
\hline
3 & 1 & 3 & 4 \\
\hline
\end{tabular}
\end{center}

\begin{enumerate}
  \item Déterminer les séries marginales de $X$ et $Y$.
  \item Déterminer les coordonnées $\overline{X}$ et $\overline{Y}$ du point moyen $G$ du nuage statistique.
  \item On admet que la variance de $X$ est $\dfrac{10}{169}$ et celle de $Y$ est $2{,}59$.
  \begin{enumerate}[label=\alph*.]
    \item Montrer que la covariance de $X$ et $Y$ est égale à $\dfrac{29}{169}$.
    \item Montrer que la droite de régression linéaire de $Y$ en $X$ est : $Y = \dfrac{29}{40}X - \dfrac{1}{20}$.
  \end{enumerate}
\end{enumerate}

\end{document}
